\documentclass[11pt,a4paper]{article}

\usepackage[utf8]{inputenc}
\usepackage[T1]{fontenc}
\usepackage[french]{babel}
\usepackage{lmodern}

\usepackage[margin=2.5cm]{geometry}
\usepackage{graphicx}
\usepackage{booktabs}
\usepackage{caption}
\usepackage{siunitx}
\sisetup{group-separator={\,}, group-minimum-digits=4, detect-all}
\usepackage{csquotes}
\usepackage[backend=biber,style=numeric]{biblatex}
\addbibresource{references.bib}
\usepackage[hidelinks]{hyperref}

\title{Cartographier les manques des ERP :\\
       une analyse des modules tiers de l'écosystème Odoo}
\author{Ameth Yade \\
        \small M2 SIGLIS --- Université de Pau et des Pays de l'Adour}
\date{\today}

\begin{document}
\maketitle

\begin{abstract}
% À écrire en dernier. Quatre à six phrases : le problème, la méthode,
% le résultat principal, la portée. C'est la seule partie que la majorité
% des lecteurs lira, et elle doit tenir seule.
\end{abstract}

\tableofcontents
\newpage

\section{Introduction}
\subsection{Le marché des ERP et le modèle modulaire}
\subsection{La question de recherche}
\subsection{Ce que ce travail apporte}

\section{Méthodologie}
\subsection{Source des données et périmètre}
\subsection{Protocole de collecte et précautions éthiques}

% Les règles pertinentes du robots.txt d'apps.odoo.com sont
% Disallow: */page/*/*, Disallow: *?*page= et Disallow: *?*order=.
% Les URL employées sont de la forme
% /apps/modules/category/<catégorie>/browse/page/<n> : la pagination est un
% segment terminal, donc hors du motif */page/*/* ; aucun paramètre page= ni
% order= n'est utilisé. Délai de 1,5 s entre requêtes, User-Agent identifié.

\subsection{Reconstitution de la liste des catégories}

% Section à part entière : la page d'accueil n'affiche que 18 catégories, dont
% la somme ne couvre que 68 \% du catalogue. 111 catégories ont été
% reconstituées par sondage, validées par un écart résiduel de
% \num{74} modules sur \num{90848}, soit \SI{0.08}{\percent}.

\subsection{Nettoyage : décisions prises et justifications}
\subsection{Table de correspondance vers un référentiel fonctionnel}
\subsection{Variables retenues}

\section{Description du corpus}
\subsection{Volumétrie}
\subsection{Répartition par fonction métier}
\subsection{Structure du marché des éditeurs}

\section{Résultats}
\subsection{Où se concentrent les modules tiers}
\subsection{Le poids des contraintes locales et réglementaires}
\subsection{La spécialisation sectorielle}
\subsection{Évolution par version : manques comblés, manques persistants}

\section{Discussion}
\subsection{Ce que révèle cette carte}
\subsection{Implications pour un projet ERP : faire ou acheter}
\subsection{Limites de l'étude}

% Rappel à tenir dans tout le document : un volume élevé de modules dans une
% fonction ne prouve pas que le standard y est mauvais. Il peut signaler un
% besoin réellement non couvert, un besoin trop spécifique pour être
% standardisé, ou une rubrique où il est simplement facile de publier.

\section{Conclusion}

\appendix
\section{Dictionnaire des données}
\section{Table de correspondance des catégories}
\section{Code source}

\printbibliography

\end{document}
